\documentclass[11pt]{article}
\usepackage[a4paper,margin=26mm]{geometry}
\usepackage{amsmath,amssymb}
\usepackage[hidelinks]{hyperref}
\setlength{\emergencystretch}{3em}
\title{Conjecture 00000004277: an uncountable maximal type chain}
\author{lizaixi01, with OpenAI Codex assistance}
\date{23 October 2026}
\begin{document}
\maketitle
\section{The necessary order-type clause}
The conjecture asserts, in particular, that every maximal chain of the
type lattice has order type $\omega+\omega^*$. This is a countable
order. We construct a genuine uncountable chain in the full standard
height-type lattice, extend it to a maximal chain, and prove that this
maximal chain cannot be order-isomorphic to that printed order.
Refuting this universal necessary clause refutes the conjunction.
We make no claim that the separate number of all maximal chains
differs from the continuum.

Here ``types'' means the classical height types of torsion-free
abelian-group theory. Bunina's fixed arXiv version~\cite[Section 3.3,
pp.~4--5]{Bunina} recalls their precise quotient, order and coordinate
operations; she calls them \emph{h-types} to distinguish them from
complete first-order logical types. We formalize this standard
combinatorial type lattice itself. We do not claim to formalize the
separate classification of rank-one groups by types.

\section{The full standard quotient and its order}
Let $\mathcal P$ be the actual set of natural primes and put
\[
 \mathcal C=(\mathbb N\cup\{\infty\})^{\mathcal P}.
\]
Lean uses all such characteristics, not only the special zero/infinity
ones used later. The standard equivalence is
\[
 a\sim b\quad\Longleftrightarrow\quad
 \{p:a_p\ne b_p\}\text{ is finite and }
 (a_p=\infty\iff b_p=\infty)\text{ for every }p.
\]
Thus changes are permitted only at finitely many \emph{finite}
coordinates. A single changed infinity coordinate already prevents
equivalence; this is not equality modulo arbitrary finite changes.

To implement the quotient, define a preorder
\[
 a\preceq b\quad\Longleftrightarrow\quad
 (a_p=\infty\Rightarrow b_p=\infty\text{ for all }p)
 \quad\text{and}\quad \{p:a_p>b_p\}\text{ is finite}.
\]
Lean proves reflexivity and transitivity, and proves that comparison
in both directions is exactly $\sim$. Its actual
\texttt{Antisymmetrization} is therefore the full type quotient.
The theorem \texttt{standard\_order} proves the exact standard bridge
\[
 [a]\le[b]\iff
 \exists a'\sim a,\ b'\sim b\quad\forall p,\ a'_p\le b'_p.
\]
For the forward implication take $a'_p=\min(a_p,b_p)$ and $b'=b$.
Only the finite bad set is changed, and the infinity guard ensures
that all changed entries are finite. The converse follows from
transitivity and the two equivalences.

\newpage
\section{Actual meet and join}
We establish a genuine \texttt{Lattice} instance on the quotient.
Its laws follow from finite unions of the bad-coordinate sets and
the infinity guards. The separately verified compatibility formulas
are
\[
 [a]\wedge[b]=[\,p\mapsto\min(a_p,b_p)\,],\qquad
 [a]\vee[b]=[\,p\mapsto\max(a_p,b_p)\,].
\]
Consequently the partial order and the operations used in the proof
are those of the standard type lattice. The proof does not replace
the full lattice by an unrelated collection of subsets.

\section{Real cuts give an uncountable chain}
Both the primes and the rationals are countably infinite. Choose a
genuine bijection $e:\mathcal P\longrightarrow\mathbb Q$. For each
real number $x$ define the full characteristic
\[
 c(x)_p=\begin{cases}
 \infty,&e(p)<x,\\
 0,&e(p)\ge x.
 \end{cases}
\]
If $x\le y$, then $c(x)\le c(y)$ coordinatewise, hence
$[c(x)]\le[c(y)]$ in the actual quotient order. If $x<y$, rational
density supplies $q$ with $x<q<y$. At $p=e^{-1}(q)$ the two
characteristics have values zero and infinity. Their infinity
patterns differ, so they are not equivalent. Lean proves that
$x\mapsto[c(x)]$ is injective and monotone.

Its actual range $C$ is therefore a chain. If $C$ were countable,
this injection into the range subtype would make $\mathbb R$
countable, a contradiction. Thus $C$ is uncountable.

\section{A maximal chain with a different order type}
Apply Mathlib's Zorn-lemma theorem
\texttt{IsChain.exists\_maxChain} to this actual chain. It produces
an actual set $M$ of height types with
\[
 \operatorname{IsMaxChain}(\le,M),\qquad C\subseteq M.
\]
Since $C$ is uncountable, so is $M$. The printed order is encoded as
the actual ordered sum $\mathbb N\mathbin{\oplus_{\rm lex}}\mathbb N^{\rm op}$,
whose underlying set is countable. Any actual order isomorphism
$M\simeq_o\mathbb N\mathbin{\oplus_{\rm lex}}\mathbb N^{\rm op}$
would therefore make $M$ countable. Lean proves its nonexistence,
and explicitly negates the universal order-type clause.

\section{Verification}
Lean 4.33.0 and Mathlib commit
\begin{center}\small\texttt{db584cd6d46c92f209a44c0f1c829460d327499d}\end{center}
are pinned. The package includes a new root build, transitive axiom
audits for all 22 theorems and the three proof-bearing instances,
and full official \texttt{leanchecker --fresh Main} replay.
No \texttt{sorry}, \texttt{native\_decide} or additional axiom is used.
OpenAI Codex assisted; a separate supervising Codex agent reviewed
source, report and PDF. This is AI review, not human certification.
Organizer acceptance is pending.
\begin{thebibliography}{1}
\bibitem{Bunina} E.~Bunina, \emph{Types in torsion free Abelian groups},
arXiv:2409.07728v1 (2024), Section~3.3, pp.~4--5.
\url{https://arxiv.org/pdf/2409.07728v1}.
\end{thebibliography}
\end{document}
